\section{Integración vertical vs. capas de modelo/aplicación}

\subsection{El reto a la percepción tradicional}

En el pasado, se creía en general que la integración vertical (modelo y producto en uno) siempre daba mejores resultados. Pero la observación real muestra que esto no es del todo cierto:

\begin{knowledgebox}{Comparación de las dos rutas}
\textbf{Integración vertical}: aplicaciones 2C como ChatGPT o Doubao (豆包), donde el modelo y el producto están estrechamente acoplados e iteran de manera conjunta.

\textbf{Arquitectura por capas}: como Claude/Gemini (la capa de modelo) más Manus (la capa de aplicación), donde la capa de modelo y la capa de aplicación requieren capacidades muy distintas y quedan a cargo, por separado, del equipo más idóneo para cada una.
\end{knowledgebox}

\begin{warningbox}{Conclusiones distintas para 2C y 2B}
En las aplicaciones 2C, la integración vertical sigue siendo válida: el producto está profundamente ligado al modelo. Pero en el extremo 2B la tendencia parece ser la contraria: a la vez que el modelo se vuelve más capaz, surgen más productos de la capa de aplicación que aprovechan bien ese modelo en distintos eslabones de la productividad.
\end{warningbox}

